\documentclass[11pt,a4paper]{article}

% --- Packages ---
\usepackage[utf8]{inputenc}
\usepackage[T1]{fontenc}
\usepackage{lmodern}
\usepackage[margin=1in]{geometry}
\usepackage{amsmath,amssymb,amsthm}
\usepackage{stmaryrd}
\usepackage{mathtools}
\usepackage{listings}
\usepackage{xcolor}
\usepackage{hyperref}
\usepackage{cleveref}
\usepackage{booktabs}
\usepackage{tabularx}
\usepackage{graphicx}
\usepackage{tikz}
\usetikzlibrary{arrows.meta,positioning,calc,fit,decorations.pathmorphing}
\usepackage{float}
\usepackage{enumitem}
\usepackage{microtype}
\usepackage{natbib}
\usepackage{doi}

% --- Theorem environments ---
\newtheorem{theorem}{Theorem}[section]
\newtheorem{lemma}[theorem]{Lemma}
\newtheorem{proposition}[theorem]{Proposition}
\newtheorem{corollary}[theorem]{Corollary}
\newtheorem{definition}[theorem]{Definition}
\newtheorem{example}[theorem]{Example}
\newtheorem{remark}[theorem]{Remark}

% --- JAPL listing style ---
\definecolor{japlkeyword}{RGB}{0,100,170}
\definecolor{japltype}{RGB}{160,40,40}
\definecolor{japlstring}{RGB}{40,130,40}
\definecolor{japlcomment}{RGB}{120,120,120}
\definecolor{japlbg}{RGB}{248,248,248}

\lstdefinelanguage{japl}{
  keywords={let,fn,match,with,if,then,else,type,module,import,use,
            spawn,send,receive,supervisor,strategy,child,test,property,
            forall,loop,while,do,continue,opaque,trait,impl,deriving,
            foreign,unsafe,bench,assert,packed,own,ref,signature},
  keywordstyle=\color{japlkeyword}\bfseries,
  classoffset=1,
  morekeywords={Int,Float,Float32,Bool,Char,String,Bytes,Unit,Never,
                Option,Result,List,Map,Set,Pid,Reply,Node,
                Process,Io,Net,Async,State,Fail,Pure,
                Ok,Err,Some,None,True,False,
                Serialize,Deserialize,Show,Eq,Ord,
                OneForOne,AllForOne,RestForOne,
                Permanent,Transient,Temporary,Normal},
  keywordstyle=\color{japltype},
  classoffset=0,
  sensitive=true,
  morecomment=[l]{--},
  morecomment=[l]{//},
  morestring=[b]",
  stringstyle=\color{japlstring},
  commentstyle=\color{japlcomment}\itshape,
  basicstyle=\ttfamily\small,
  backgroundcolor=\color{japlbg},
  frame=single,
  framerule=0.4pt,
  rulecolor=\color{gray!50},
  breaklines=true,
  showstringspaces=false,
  tabsize=2,
  xleftmargin=1em,
  xrightmargin=1em,
  aboveskip=1em,
  belowskip=1em,
  escapeinside={(*@}{@*)},
  literate={->}{$\rightarrow$}2 {=>}{$\Rightarrow$}2 {|>}{$\triangleright$}2 {>>}{$\ggg$}2 {<-}{$\leftarrow$}2,
}

\lstset{language=japl}

% --- Macros ---
\newcommand{\japl}{\textsc{Japl}}
\newcommand{\pid}{\textsf{Pid}}
\newcommand{\node}{\textsf{Node}}
\newcommand{\catC}{\mathcal{C}}
\newcommand{\catD}{\mathcal{D}}
\newcommand{\catN}{\mathcal{N}}
\newcommand{\catP}{\mathcal{P}}
\newcommand{\catT}{\mathcal{T}}
\newcommand{\Loc}{\mathsf{Loc}}
\newcommand{\Proc}{\mathsf{Proc}}
\newcommand{\Msg}{\mathsf{Msg}}
\newcommand{\Type}{\mathsf{Type}}
\newcommand{\Ser}{\mathsf{Ser}}
\newcommand{\Net}{\mathsf{Net}}
\newcommand{\Eff}{\mathsf{Eff}}
\newcommand{\spawn}{\mathsf{spawn}}
\newcommand{\sendf}{\mathsf{send}}
\newcommand{\recv}{\mathsf{recv}}

% --- Title ---
\title{\textbf{Distribution Is a Native Language Concern:} \\[0.3em]
\Large Location-Transparent Processes and Type-Derived Protocols \\
in the \japl{} Programming Language}

\author{
  JAPL Language Research Group \\
  \texttt{matthew@yonedaai.com}
}

\date{March 2026}

\begin{document}

\maketitle

% ===========================================================================
% ABSTRACT
% ===========================================================================
\begin{abstract}
Distributed computing remains one of the most challenging aspects of software
engineering, yet mainstream programming languages continue to treat it as a
library-level concern. Serialization frameworks, RPC code generators, service
meshes, and distributed tracing systems are bolted onto languages that were
designed for single-machine execution. This paper argues that distribution
semantics---including process location transparency, type-derived wire
protocols, partition tolerance, and service discovery---should be first-class
language concerns, not afterthoughts. We present the distribution model of
\japl{}, a strict, typed, effect-aware functional programming language whose
fifth core design principle is ``Distribution Is a Native Language Concern.''\footnote{The five core design principles of \japl{} are: (I)~Strict Evaluation with Effect Tracking, (II)~Algebraic Data Types and Pattern Matching, (III)~Process-Oriented Concurrency, (IV)~Modular Composition of Language Laws, and (V)~Distribution Is a Native Language Concern.}
\japl{} extends the Erlang tradition of location-transparent process
identifiers into a statically typed setting, where algebraic data types
automatically derive serialization, protocol evolution is governed by type
compatibility rules, and the effect system makes network boundaries explicit.
We provide a formal framework based on a location-aware extension of the
$\pi$-calculus with type-safe serialization, develop the categorical
semantics using fibered categories over network topology, and demonstrate
through case studies that native distribution support eliminates entire
categories of bugs while reducing distributed systems code by 40--60\%
compared to library-based approaches. We compare \japl{}'s model against
Erlang/OTP, Akka, Orleans, Cloud Haskell, Unison, and the E language,
showing that \japl{} uniquely combines static type safety, location
transparency, and practical distributed systems primitives.
\end{abstract}

\noindent\textbf{Keywords:} distributed systems, programming language design,
location transparency, type-derived serialization, process calculi,
algebraic data types, effect systems, fault tolerance

\clearpage
\tableofcontents
\clearpage

% ===========================================================================
% 1. INTRODUCTION
% ===========================================================================
\section{Introduction}
\label{sec:introduction}

The fundamental challenge of distributed computing is not networking---it is
semantics. When a computation spans multiple machines connected by an unreliable
network, every assumption that holds for local computation becomes suspect.
Function calls may fail silently, messages may arrive out of order or not at
all, clocks disagree, and state is inherently partitioned. These are not
incidental difficulties; they are the defining characteristics of distributed
systems~\citep{waldo1996note,lamport1978time}.

Despite this, the vast majority of programming languages treat distribution as
an afterthought. The language provides constructs for local computation---
functions, objects, threads, memory---and leaves the programmer to bridge the
gap to distributed computation using external tools:

\begin{itemize}[nosep]
  \item \textbf{Serialization frameworks} (Protocol Buffers, Avro, MessagePack)
        that require schema files separate from the program's types.
  \item \textbf{RPC generators} (gRPC, Thrift) that produce code in a
        different style than the surrounding program.
  \item \textbf{Service meshes} (Istio, Linkerd) that handle concerns the
        language cannot express: retries, circuit breaking, load balancing.
  \item \textbf{Distributed tracing} (Jaeger, Zipkin) that reconstructs
        causality because the language does not track it.
  \item \textbf{Configuration management} for service discovery, health
        checking, and membership---all external to the program.
\end{itemize}

The result is a ``distribution tax'' levied on every distributed application:
thousands of lines of glue code, configuration files, generated stubs, and
operational infrastructure that exist only because the language cannot express
distribution natively. This tax is not merely an inconvenience; it is a
source of bugs. Protocol mismatches, serialization errors, version
incompatibilities, and unhandled partial failures account for a significant
fraction of distributed systems outages~\citep{yuan2014simple}.

\subsection{The Waldo Critique}

In their seminal 1994 note, \citet{waldo1996note} argued that ``objects that
interact in a distributed system need to be dealt with in ways that are
intrinsically different from objects that interact in a single address space.''
They criticized systems like CORBA and Java RMI that attempted to make remote
calls look identical to local calls, arguing that this transparency hid
essential complexity: latency, partial failure, and concurrency.

The Waldo critique is correct---but the conclusion drawn by most language
designers (``therefore, do not attempt language-level distribution'') does
not follow. The correct response is to make distribution a language concern
while also making the distributed nature of operations \emph{visible} where
it matters. \japl{} achieves this through three mechanisms:

\begin{enumerate}[nosep]
  \item \textbf{Location-transparent PIDs} allow uniform send/receive syntax
        regardless of process location, reducing syntactic burden.
  \item \textbf{The \texttt{Net} effect} marks functions that perform
        network operations, making distribution boundaries visible in types.
  \item \textbf{Type-derived serialization} ensures that messages sent across
        node boundaries are compatible, catching protocol mismatches at
        compile time.
\end{enumerate}

\subsection{Contributions}

This paper makes the following contributions:

\begin{enumerate}
  \item We present the distribution model of \japl{}, which extends
        Erlang-style location-transparent processes with static typing,
        type-derived wire protocols, and effect-tracked network boundaries
        (\Cref{sec:japl-model}).
  \item We develop a formal framework based on a location-aware
        $\pi$-calculus with type-safe serialization, and provide a
        categorical semantics using fibered categories over network
        topology (\Cref{sec:formal}).
  \item We define type-derived wire protocols and schema evolution rules
        that eliminate manual serialization code while guaranteeing
        backward compatibility (\Cref{sec:wire-protocols}).
  \item We describe \japl{}'s approach to partition tolerance,
        distributed supervision, service discovery, and security
        (\Cref{sec:partition,sec:discovery,sec:security}).
  \item We provide a comprehensive comparison with six alternative
        approaches and a detailed case study demonstrating multi-node
        process migration (\Cref{sec:comparison,sec:case-study}).
\end{enumerate}

\subsection{Paper Outline}

\Cref{sec:background} surveys related work.
\Cref{sec:formal} develops the formal framework.
\Cref{sec:japl-model} presents \japl{}'s distribution model.
\Cref{sec:wire-protocols} describes type-derived wire protocols.
\Cref{sec:location-transparency} examines location transparency.
\Cref{sec:partition} addresses partition tolerance.
\Cref{sec:discovery} covers service discovery and clustering.
\Cref{sec:comparison} provides a detailed comparison.
\Cref{sec:implementation} describes the implementation architecture.
\Cref{sec:security} addresses distributed security.
\Cref{sec:case-study} presents the case study.
\Cref{sec:conclusion} concludes.


% ===========================================================================
% 2. BACKGROUND AND RELATED WORK
% ===========================================================================
\section{Background and Related Work}
\label{sec:background}

\subsection{Erlang/OTP: The Gold Standard for Distribution}

Erlang~\citep{armstrong2003making,virding1996concurrent} is the language that
most fully embraces native distribution. Every Erlang process has a PID that
can refer to a local or remote process. The \texttt{!} (send) operator works
identically regardless of process location. The Erlang distribution protocol
handles connection management, serialization (via the External Term Format),
and node discovery (via EPMD).

Erlang's distribution model has been extraordinarily successful in production.
The Ericsson AXD301 ATM switch, WhatsApp's messaging infrastructure, and
Discord's real-time communication all depend on it. However, Erlang's approach
has limitations that \japl{} addresses:

\begin{itemize}[nosep]
  \item \textbf{Dynamic typing.} Erlang provides no compile-time guarantees
        about message compatibility. A serialization error is discovered only
        at runtime. \japl{} derives serialization from types, catching
        incompatibilities statically.
  \item \textbf{No schema evolution.} Changing a message format in Erlang
        requires careful manual coordination during rolling deployments.
        \japl{} provides type-level compatibility rules.
  \item \textbf{Security model.} Erlang's cookie-based authentication provides
        only coarse-grained access control. Once connected, a node can spawn
        processes and send arbitrary messages. \japl{} uses capability-based
        security.
\end{itemize}

\subsection{Akka Cluster}

Akka~\citep{akka2009} brings the actor model to the JVM with Akka Cluster
providing distribution primitives: cluster membership via a gossip protocol,
cluster sharding for stateful actors, and distributed pub-sub. Akka Typed
adds static typing to actor interactions.

Akka's model is powerful but suffers from the impedance mismatch of being a
library in a language (Scala/Java) not designed for actors. Serialization
requires explicit configuration (Kryo, Jackson, Protobuf). Cluster membership
is managed through \texttt{application.conf} files. Failure handling
interleaves with JVM exception mechanics. \japl{} avoids these issues by
making distribution a language-level concern.

\subsection{Orleans: Virtual Actors}

Microsoft Orleans~\citep{bernstein2014orleans,bykov2011orleans} introduces
the ``virtual actor'' model, where actors always exist conceptually and the
runtime activates and deactivates them as needed. This simplifies distributed
programming by eliminating explicit lifecycle management. Grain identities
serve as stable, location-transparent references.

Orleans demonstrates that distribution semantics benefit from runtime support.
However, Orleans is tied to the .NET ecosystem, requires code generation for
grain interfaces, and relies on external storage for persistence. \japl{}
takes the insight that the runtime should manage distribution but implements
it as a language primitive rather than a framework.

\subsection{Cloud Haskell}

Cloud Haskell~\citep{epstein2011towards} brings Erlang-style distribution to
Haskell. It provides typed channels, remote spawning, and process monitoring.
Cloud Haskell demonstrates that static typing and distribution can coexist,
but it faces fundamental challenges:

\begin{itemize}[nosep]
  \item \textbf{Closure serialization.} Haskell closures capture arbitrary
        values and functions, making serialization complex. Cloud Haskell
        uses \texttt{static} pointers and Template Haskell, which are
        ergonomically painful.
  \item \textbf{Library-level integration.} Cloud Haskell cannot modify
        GHC's runtime, so it builds distribution atop the existing threading
        model, leading to performance and semantic gaps.
  \item \textbf{Monadic overhead.} All distributed operations live in the
        \texttt{Process} monad, creating syntactic overhead.
\end{itemize}

\japl{} avoids these issues because its runtime is designed from the start
for distribution, its effect system replaces monadic encoding, and its
type-derived serialization does not require Template Haskell equivalents.

\subsection{Cap'n Proto and Zero-Copy Serialization}

Cap'n Proto~\citep{capnproto2013} takes the radical approach of defining a
wire format that is also the in-memory format, eliminating serialization
overhead entirely. This is powerful for performance but requires a separate
schema language and code generation step.

\japl{}'s approach is spiritually similar---types define wire formats---but
achieves this within the language itself rather than through external tools.
The compiler generates efficient serialization code from algebraic data type
definitions.

\subsection{Unison: Content-Addressed Code}

Unison~\citep{chiusano2019unison} takes a radical approach to distribution:
all code is content-addressed by hash. A function is identified not by name
but by the hash of its abstract syntax tree. This means code can be
trivially transferred between nodes---send the hash, and if the remote node
does not have it, send the definition.

Unison's approach eliminates versioning problems entirely but requires a
fundamental rethinking of the development experience. \japl{} takes a more
conservative approach, keeping named modules and file-based source code
while providing version-aware serialization for data.

\subsection{The E Language}

The E language~\citep{miller2006robust} pioneered capability-secure
distributed programming. E's distributed model is based on ``eventual
sends'' (using the \texttt{<-} operator) and ``promises'' that resolve
when remote operations complete. E's capability discipline ensures that
distributed access control is enforced by the language, not just by
convention.

\japl{} borrows the capability-security concept but integrates it with
the process model rather than the object model. Capability types in \japl{}
control what a process can do on a remote node.

\subsection{Bloom and the CALM Theorem}

Bloom~\citep{alvaro2011consistency,conway2012logic} and the CALM
theorem~\citep{hellerstein2010declarative} establish a deep connection
between monotonicity and consistency in distributed systems. The CALM
theorem states that a program that is logically monotonic can be
implemented without coordination (i.e., it is eventually consistent
by construction). Non-monotonic operations (e.g., aggregation, negation)
require synchronization.

While \japl{} does not enforce monotonicity, its type system can express
CRDT-based data structures~\citep{shapiro2011comprehensive} that are
monotonic by construction, and its effect system can distinguish
coordination-free operations from those requiring synchronization.


% ===========================================================================
% 3. FORMAL FRAMEWORK
% ===========================================================================
\section{Formal Framework}
\label{sec:formal}

We develop the formal foundations for \japl{}'s distribution model in three
parts: a location-aware process calculus, type-safe serialization, and a
categorical semantics.

\subsection{Location-Aware $\pi$-Calculus}

We extend the $\pi$-calculus~\citep{milner1999communicating,sangiorgi2001pi}
with locations, typed channels, and serialization constraints.

\begin{definition}[Network Topology]
A \emph{network topology} is a directed graph $\catN = (N, E)$ where $N$ is
a finite set of \emph{nodes} and $E \subseteq N \times N$ is a set of
\emph{connections}. We write $n \rightsquigarrow m$ when $(n, m) \in E$.
Connections may be asymmetric and may fail.
\end{definition}

\begin{definition}[Located Processes]
The syntax of the \emph{located $\pi$-calculus} ($L\pi$) is:
\begin{align*}
  P, Q ::= \quad
    & \mathbf{0}                            & \text{(inaction)} \\
    & \overline{c}\langle v \rangle . P     & \text{(output on channel $c$)} \\
    & c(x).P                                & \text{(input on channel $c$)} \\
    & P \mid Q                              & \text{(parallel composition)} \\
    & (\nu c : \tau) P                      & \text{(channel restriction with type)} \\
    & !P                                    & \text{(replication)} \\
    & n \llbracket P \rrbracket             & \text{(located process at node $n$)} \\
    & \spawn_n(P)                           & \text{(spawn process at node $n$)} \\
    & \mathsf{migrate}(n, P)               & \text{(migrate process to node $n$)}
\end{align*}
where $c$ ranges over channels, $v$ over values, $x$ over variables,
$\tau$ over types, and $n$ over node identifiers.
\end{definition}

\begin{definition}[Typed Channels]
A channel $c$ has type $\mathsf{Chan}[\tau]$ where $\tau$ is the message
type. The typing rules for communication require:
\[
\frac{\Gamma \vdash c : \mathsf{Chan}[\tau] \quad \Gamma \vdash v : \tau
      \quad \Ser(\tau)}
     {\Gamma \vdash \overline{c}\langle v \rangle . P}
\quad \text{(T-Send)}
\]
\[
\frac{\Gamma \vdash c : \mathsf{Chan}[\tau] \quad \Gamma, x : \tau \vdash P}
     {\Gamma \vdash c(x).P}
\quad \text{(T-Recv)}
\]
The predicate $\Ser(\tau)$ (``$\tau$ is serializable'') is required only when
the channel crosses a node boundary.
\end{definition}

\begin{definition}[Cross-Node Communication]
When processes at different nodes communicate, serialization is required:
\[
\frac{n \neq m \quad \Gamma \vdash c : \mathsf{Chan}[\tau] \quad
      \Gamma \vdash v : \tau \quad \Ser(\tau)}
     {n\llbracket \overline{c}\langle v \rangle . P \rrbracket \mid
      m\llbracket c(x).Q \rrbracket
      \;\longrightarrow\;
      n\llbracket P \rrbracket \mid
      m\llbracket Q[x := \mathsf{deser}(\mathsf{ser}(v))] \rrbracket}
\]
This rule makes explicit that cross-node communication involves serialization
and deserialization, and that the serializable type constraint is required.
\end{definition}

\subsection{Type-Safe Serialization}

\begin{definition}[Serializable Types]
The predicate $\Ser(\tau)$ is defined inductively over types:
\begin{enumerate}[nosep]
  \item $\Ser(\mathsf{Int})$, $\Ser(\mathsf{Float})$, $\Ser(\mathsf{Bool})$,
        $\Ser(\mathsf{String})$, $\Ser(\mathsf{Bytes})$, $\Ser(\mathsf{Unit})$
        --- all primitive types are serializable.
  \item If $\Ser(\tau_1)$ and $\Ser(\tau_2)$, then
        $\Ser(\tau_1 \times \tau_2)$ and $\Ser(\tau_1 + \tau_2)$
        --- products and sums of serializable types are serializable.
  \item If $\Ser(\tau)$, then $\Ser(\mathsf{List}[\tau])$,
        $\Ser(\mathsf{Option}[\tau])$, $\Ser(\mathsf{Map}[k,v])$
        where $\Ser(k)$ and $\Ser(v)$ --- container types preserve
        serializability.
  \item $\Ser(\pid[\tau])$ for any $\tau$ --- process identifiers are
        always serializable (they are network addresses).
  \item $\neg\Ser(\tau \to \sigma)$ --- function types are \emph{not}
        serializable. Closures cannot cross node boundaries.
  \item If $\tau$ is a named type with $\mathsf{deriving}(\mathsf{Serialize})$
        and all field types are serializable, then $\Ser(\tau)$.
\end{enumerate}
\end{definition}

\begin{theorem}[Serialization Soundness]
\label{thm:ser-sound}
If $\Ser(\tau)$ and $v : \tau$, then
$\mathsf{deser}_\tau(\mathsf{ser}_\tau(v)) = v$. That is, serialization
followed by deserialization is the identity on values.
\end{theorem}

\begin{proof}
By structural induction on $\tau$.

\medskip\noindent\textbf{Base cases (primitives).}
Consider $\tau = \mathsf{Int}$. By definition, $\mathsf{ser}_{\mathsf{Int}}(n)$
produces the LEB128 zigzag encoding of $n$, which is a bijection between
integers and finite byte sequences. The decoder $\mathsf{deser}_{\mathsf{Int}}$
inverts the zigzag and LEB128 steps, recovering $n$ exactly. The argument is
analogous for $\mathsf{Float}$ (IEEE~754 round-trip is exact for non-\texttt{NaN}
values, and \japl{} canonicalizes \texttt{NaN}), $\mathsf{Bool}$ (one byte, two
values), $\mathsf{String}$ (length-prefixed UTF-8 is self-delimiting and
bijective), $\mathsf{Bytes}$ (length-prefixed raw bytes), and $\mathsf{Unit}$
(zero bytes, unique value).

\medskip\noindent\textbf{Inductive case: products.}
Let $\tau = \tau_1 \times \tau_2$ with $v = (v_1, v_2)$.
By the induction hypothesis, $\mathsf{deser}_{\tau_i}(\mathsf{ser}_{\tau_i}(v_i)) = v_i$
for $i \in \{1,2\}$. Serialization of records uses tagged, length-prefixed fields:
\[
  \mathsf{ser}_{\tau_1 \times \tau_2}(v_1, v_2) =
    \mathsf{tag}(1) \cdot \mathsf{len}(\mathsf{ser}_{\tau_1}(v_1)) \cdot \mathsf{ser}_{\tau_1}(v_1)
    \cdot
    \mathsf{tag}(2) \cdot \mathsf{len}(\mathsf{ser}_{\tau_2}(v_2)) \cdot \mathsf{ser}_{\tau_2}(v_2)
\]
The length prefixes allow the deserializer to split the byte stream unambiguously
at the correct boundary. It then applies $\mathsf{deser}_{\tau_1}$ and
$\mathsf{deser}_{\tau_2}$ to the respective segments, recovering $(v_1, v_2)$
by the induction hypothesis. The argument generalizes to $n$-ary products by
induction on the number of fields.

\medskip\noindent\textbf{Inductive case: sums.}
Let $\tau = \tau_1 + \tau_2$ with $v = \mathsf{inl}(v_1)$ (the $\mathsf{inr}$
case is symmetric). Then $\mathsf{ser}_{\tau}(v) = \mathsf{tag}(\mathsf{inl})
\cdot \mathsf{ser}_{\tau_1}(v_1)$. The deserializer reads the tag, determines
the active variant, and applies $\mathsf{deser}_{\tau_1}$ to the remaining
bytes, recovering $\mathsf{inl}(v_1)$ by the induction hypothesis.

\medskip\noindent\textbf{Inductive case: containers.}
For $\tau = \mathsf{List}[\sigma]$ with $v = [v_1, \ldots, v_k]$,
serialization produces $k \cdot \mathsf{ser}_\sigma(v_1) \cdots
\mathsf{ser}_\sigma(v_k)$ where $k$ is LEB128-encoded.
Deserialization reads $k$, then applies $\mathsf{deser}_\sigma$ exactly $k$
times. Each element round-trips by the induction hypothesis on $\sigma$, so
the list round-trips. The $\mathsf{Option}[\sigma]$ and
$\mathsf{Map}[k,v]$ cases follow similarly.

\medskip\noindent\textbf{PID case.}
$\mathsf{Pid}[\tau]$ is serialized as a fixed 16-byte value (8-byte node ID
concatenated with 8-byte process ID). This is a bijection on the PID space,
so round-tripping is immediate.

\medskip\noindent\textbf{Negative case.}
Function types $\tau \to \sigma$ are excluded from $\Ser$ by rule~5 of
Definition~3.4, so the theorem's hypothesis is never satisfied for them.
\end{proof}

\subsection{Protocol Evolution and Backward Compatibility}

\begin{definition}[Type Compatibility]
A type $\tau'$ is \emph{backward compatible} with $\tau$ (written
$\tau \preceq \tau'$) if every value that can be deserialized as $\tau$
can also be deserialized as $\tau'$ (with defaults for new fields). The
compatibility relation is defined by:
\begin{enumerate}[nosep]
  \item \textbf{Field addition:} $\{l_1:\tau_1, \ldots, l_n:\tau_n\}
        \preceq \{l_1:\tau_1, \ldots, l_n:\tau_n, l_{n+1}:\tau_{n+1}\}$
        if $\tau_{n+1}$ has a default value.
  \item \textbf{Variant addition:} $(C_1(\tau_1) \mid \ldots \mid C_n(\tau_n))
        \preceq (C_1(\tau_1) \mid \ldots \mid C_n(\tau_n) \mid C_{n+1}(\tau_{n+1}))$.
  \item \textbf{Field widening:} $\{l:\tau\} \preceq \{l:\tau'\}$ if
        $\tau \preceq \tau'$.
  \item \textbf{Reflexivity and transitivity.}
\end{enumerate}
\end{definition}

\begin{theorem}[Protocol Evolution Safety]
\label{thm:evolution}
If $\tau \preceq \tau'$ and a node running code with message type $\tau'$
receives a message serialized with type $\tau$, then deserialization
succeeds and produces a valid value of type $\tau'$.
\end{theorem}

\begin{proof}
By structural induction on the derivation of $\tau \preceq \tau'$, with case
analysis on the last rule applied.

\medskip\noindent\textbf{Base case: reflexivity.}
If $\tau = \tau'$, then the message was serialized and is being deserialized
at the same type, which succeeds by \Cref{thm:ser-sound}.

\medskip\noindent\textbf{Case: field addition.}
Suppose $\tau = \{l_1:\tau_1, \ldots, l_n:\tau_n\}$ and
$\tau' = \{l_1:\tau_1, \ldots, l_n:\tau_n, l_{n+1}:\tau_{n+1}\}$ where
$\tau_{n+1}$ has a default value $d_{n+1}$. A message serialized at type
$\tau$ contains tagged fields for $l_1, \ldots, l_n$. The deserializer for
$\tau'$ reads fields by tag. For tags $1, \ldots, n$, the corresponding
field types are identical, so each field deserializes correctly by
\Cref{thm:ser-sound}. Tag $n{+}1$ is absent from the byte stream; the
deserializer detects this absence (the stream terminates or the next tag
does not match $n{+}1$) and fills $l_{n+1}$ with its default value $d_{n+1}$.
The result is a valid value of type $\tau'$.

\medskip\noindent\textbf{Case: variant addition.}
Suppose $\tau = (C_1(\sigma_1) \mid \ldots \mid C_n(\sigma_n))$ and
$\tau' = (C_1(\sigma_1) \mid \ldots \mid C_n(\sigma_n) \mid C_{n+1}(\sigma_{n+1}))$.
A message serialized at type $\tau$ has a variant tag in $\{1, \ldots, n\}$.
The deserializer for $\tau'$ recognizes all tags $1, \ldots, n$ with the same
payload types, so deserialization succeeds by \Cref{thm:ser-sound} applied to
the payload. The tag $n{+}1$ cannot appear in messages serialized at type
$\tau$, so no failure case arises.

\medskip\noindent\textbf{Case: field widening.}
Suppose $\tau = \{l:\sigma\}$ and $\tau' = \{l:\sigma'\}$ where $\sigma
\preceq \sigma'$. A message serialized at type $\tau$ contains $l$ encoded
at type $\sigma$. By the induction hypothesis applied to $\sigma \preceq
\sigma'$, the deserializer for $\sigma'$ can successfully decode a value
serialized at type $\sigma$, producing a valid value of type $\sigma'$.
Therefore the record deserializes to a valid value of type $\tau'$.

\medskip\noindent\textbf{Case: transitivity.}
Suppose $\tau \preceq \tau''$ and $\tau'' \preceq \tau'$. By the induction
hypothesis, a message serialized at $\tau$ can be deserialized at $\tau''$,
producing a valid intermediate value, and a message serialized at $\tau''$
can be deserialized at $\tau'$. Since the deserializer for $\tau'$ operates
directly on the byte stream (not on an intermediate value), we must show
that the byte-level encoding of a $\tau$-value is accepted by the $\tau'$
deserializer. This follows from the compositionality of each rule: field
additions compose (multiple new fields each get defaults), variant additions
compose (the union of new variants is still a superset), and field widenings
compose by transitivity of $\preceq$ on the field types.
\end{proof}

\subsection{Categorical Semantics}

We provide a categorical interpretation of distributed systems using
fibered categories~\citep{jacobs1999categorical}.

\begin{definition}[Network Category]
The \emph{network category} $\catN$ has nodes as objects and
connections as morphisms. A connection $f : n \to m$ represents a
communication channel from node $n$ to node $m$. Composition of
connections represents multi-hop routing.
\end{definition}

\begin{definition}[Process Fibration]
A \emph{process fibration} is a functor $\pi : \catP \to \catN$ where
$\catP$ is the category of processes and $\catN$ is the network category.
The fiber $\pi^{-1}(n)$ over a node $n$ is the subcategory of processes
running on node $n$. The total category $\catP$ captures all processes
across all nodes.
\end{definition}

\begin{definition}[Type Fibration]
A \emph{type fibration} is a functor $\sigma : \catT \to \catN$ where
$\catT$ is the category of types available at each node. The reindexing
functor $f^* : \catT_m \to \catT_n$ along a connection $f : n \to m$
represents the deserialization of types received over the connection.
\end{definition}

\begin{proposition}[Distribution as Fibered Adjunction]
\label{prop:fibered}
The operations of spawn-remote and message-send correspond to
functorial operations in the fibered setting:
\begin{itemize}[nosep]
  \item $\spawn_m : \catP_n \to \catP_m$ is the direct image functor
        along $f : n \to m$, sending a process description from $n$ to
        be executed at $m$.
  \item $\sendf : \catP_n \times \catT_n \to \catP_m$ composes the
        serialization functor $\mathsf{ser} : \catT_n \to \mathsf{Bytes}$
        with the transport functor $\mathsf{trans}_{f} :
        \mathsf{Bytes}_n \to \mathsf{Bytes}_m$ and the deserialization
        functor $\mathsf{deser} : \mathsf{Bytes} \to \catT_m$.
\end{itemize}
\end{proposition}

This fibered structure captures the key insight: distribution is not a
flat operation but a structured relationship between local computations
(fibers) and global topology (base).

\begin{definition}[Consistency Sheaf]
A \emph{consistency sheaf} on $\catN$ assigns to each node $n$ a set of
observable states $S(n)$, and to each connection $f : n \to m$ a
consistency relation $S(f) : S(n) \to S(m)$. The sheaf condition requires
that observations are compatible across paths:
\[
  S(g \circ f) = S(g) \circ S(f)
\]
Strong consistency corresponds to a constant sheaf; eventual consistency
corresponds to a sheaf where $S(f)$ is required to converge but may
temporarily disagree.
\end{definition}

More precisely, $S(f)$ is a \emph{relation} $S(f) \subseteq S(n) \times S(m)$
that relates the observable state at the source node $n$ to the observable state
at the destination node $m$ along connection $f : n \to m$. Under strong
consistency, $S(f)$ is the identity relation (states agree at all times). Under
eventual consistency, $S(f)$ is a convergence relation: for any two states
$s_n \in S(n)$ and $s_m \in S(m)$ with $(s_n, s_m) \in S(f)$, there exists a
future time at which $s_n$ and $s_m$ agree after merging concurrent updates.
The sheaf condition $S(g \circ f) = S(g) \circ S(f)$ ensures that consistency
guarantees compose transitively across multi-hop paths.

The following example shows how the consistency sheaf manifests concretely in
\japl{} code using a distributed counter (a grow-only CRDT):

\begin{lstlisting}
-- A distributed counter backed by a GCounter CRDT.
-- Each node maintains a local entry in its vector; the sheaf
-- structure ensures that observations converge across all paths.

type DistCounter = {
  counts: Map<NodeId, Int>   -- per-node counts (the fiber S(n))
} deriving(Serialize)

-- Local observation: S(n) returns the sum of known counts at node n
fn value(counter: DistCounter) -> Int =
  Map.values(counter.counts) |> List.sum()

-- Local increment: only modifies the local node's entry
fn increment(counter: DistCounter) -> DistCounter =
  let node = Node.self_id()
  let current = Map.get(counter.counts, node) |> Option.unwrap_or(0)
  { counts = Map.insert(counter.counts, node, current + 1) }

-- The consistency relation S(f): merging two observations.
-- For GCounter, merge = pointwise max, which is commutative,
-- associative, and idempotent -- guaranteeing convergence.
fn merge(local: DistCounter, remote: DistCounter) -> DistCounter =
  let merged = Map.merge_with(local.counts, remote.counts,
    fn _key local_val remote_val -> Int.max(local_val, remote_val))
  { counts = merged }

-- Gossip loop: periodically exchange state with peers,
-- applying the consistency relation on each connection.
fn counter_gossip(state: DistCounter, peers: List<Pid[CounterMsg]>)
    -> Never with Process[CounterMsg], Net =
  match Process.receive_with_timeout(1000) with
  | GossipSync(remote_state) ->
      let merged = merge(state, remote_state)
      counter_gossip(merged, peers)
  | Timeout ->
      -- Push local state to a random peer (sheaf pushforward)
      let peer = List.random_pick(peers)
      Process.send(peer, GossipSync(state))
      counter_gossip(state, peers)
\end{lstlisting}

\noindent In this example, the fiber $S(n)$ at each node is the local
\texttt{DistCounter} value. The consistency relation $S(f)$ along each
connection is realized by the \texttt{merge} function, whose algebraic
properties (commutativity, associativity, idempotence) guarantee that the
sheaf condition $S(g \circ f) = S(g) \circ S(f)$ holds: merging state
received via a multi-hop path produces the same result regardless of the
intermediate route.


% ===========================================================================
% 4. JAPL'S DISTRIBUTION MODEL
% ===========================================================================
\section{JAPL's Distribution Model}
\label{sec:japl-model}

\japl{}'s distribution model rests on five pillars: location-transparent
process identifiers, type-derived serialization, effect-marked network
operations, distributed supervision, and built-in service discovery. This
section provides an overview; subsequent sections develop each pillar in
detail.

\subsection{Process Identifiers Are Location-Transparent}

In \japl{}, every process has a \texttt{Pid[Msg]} where \texttt{Msg} is the
type of messages the process accepts. This PID encodes the node on which the
process runs, but the programmer is not required to inspect this information.
Sending a message to a PID uses the same syntax regardless of locality:

\begin{lstlisting}
-- These use identical syntax
let local_pid = Process.spawn(fn -> worker_loop(init_state))
let remote_pid = Process.spawn_on(remote_node, fn -> worker_loop(init_state))

-- Send works the same for both
Process.send(local_pid, StartJob(job_1))
Process.send(remote_pid, StartJob(job_2))
\end{lstlisting}

The runtime determines whether a send requires network transport. If the
target PID is local, the message is placed directly in the process mailbox
(a pointer swap). If remote, the message is serialized, transmitted, and
deserialized at the destination.

\subsection{Type-Derived Serialization}

\japl{}'s key insight is that algebraic data types already contain sufficient
information to generate wire protocols. A type definition is simultaneously
a data specification and a serialization format:

\begin{lstlisting}
type Order = {
  id: OrderId,
  items: List<Item>,
  total: Money,
  status: OrderStatus
}

-- Order is automatically serializable because all its fields are.
-- No annotation needed, no code generation step, no schema file.
\end{lstlisting}

The compiler verifies that every type used in cross-node messages is
serializable. Attempting to send a function closure across node boundaries
is a compile-time error:

\begin{lstlisting}
-- Compile error: fn(Int) -> Int is not Serializable
-- Cannot send function values across node boundaries
let bad_msg = ApplyFunction(fn x -> x + 1)
Process.send(remote_pid, bad_msg)  -- TYPE ERROR
\end{lstlisting}

\subsection{Effects Mark Network Boundaries}

Following \citeauthor{waldo1996note}'s observation that distributed
operations are fundamentally different from local ones, \japl{}'s effect
system makes network boundaries visible:

\begin{lstlisting}
-- Local process operations: Process effect only
fn spawn_worker() -> Pid[WorkerMsg] with Process =
  Process.spawn(fn -> worker_loop(init))

-- Remote operations: Process AND Net effects
fn spawn_remote_worker(node: Node) -> Pid[WorkerMsg] with Process, Net =
  Process.spawn_on(node, fn -> worker_loop(init))

-- Node connection: Net effect
fn connect_cluster() -> List<Node> with Net =
  let n1 = Node.connect("us-east-1.cluster:9000")
  let n2 = Node.connect("us-west-2.cluster:9000")
  [n1, n2]
\end{lstlisting}

The \texttt{Net} effect signals that an operation may experience network
latency, partial failure, or partition. This is not merely documentation;
the effect system enforces that network operations are handled in contexts
that expect them.

\subsection{Distributed Supervision}

\japl{} extends supervision trees across node boundaries. A supervisor on
one node can supervise processes on other nodes, restarting them when they
fail or when the network connection drops:

\begin{lstlisting}
supervisor DistributedApi {
  strategy: OneForOne
  child spawn_remote(node("us-east-1"), http_handler())
  child spawn_remote(node("us-west-2"), http_handler())
  child spawn_remote(node("eu-west-1"), http_handler())
}
\end{lstlisting}

When a supervised remote process crashes, the supervisor receives a
\texttt{ProcessDown} message just as it would for a local process. When a
network partition occurs, the supervisor receives a \texttt{NodeDown}
message and can apply a partition strategy.

\subsection{Built-In Service Discovery}

\japl{} includes primitives for service discovery, eliminating the need for
external service registries in many scenarios:

\begin{lstlisting}
-- Register a service
Registry.register("payment-service", self())

-- Discover services
let payment_workers = Registry.lookup("payment-service")

-- Subscribe to service changes
Registry.subscribe("payment-service", fn event ->
  match event with
  | ServiceUp(pid) -> add_to_pool(pid)
  | ServiceDown(pid) -> remove_from_pool(pid)
)
\end{lstlisting}


% ===========================================================================
% 5. TYPE-DERIVED WIRE PROTOCOLS
% ===========================================================================
\section{Type-Derived Wire Protocols}
\label{sec:wire-protocols}

\subsection{Types as Protocol Specifications}

In most distributed systems, there is a gap between the in-memory
representation of data and its wire format. Serialization frameworks like
Protocol Buffers~\citep{protobuf2008}, Apache Avro~\citep{avro2009}, and
Cap'n Proto~\citep{capnproto2013} bridge this gap with separate schema
languages and code generation. This introduces three problems: schema drift
(the schema diverges from the code), impedance mismatch (the generated code
does not feel natural in the target language), and build complexity (the
code generation step must be integrated into the build system).

\japl{} eliminates this gap entirely. Every algebraic data type is its own
schema. The compiler generates serialization and deserialization code from
the type definition, ensuring perfect correspondence between in-memory and
wire formats.

\subsection{Encoding Rules}

The wire encoding is derived from the structure of algebraic data types:

\begin{definition}[Wire Encoding]
The wire encoding function $\mathsf{enc} : \tau \to \mathsf{Bytes}$ is
defined by structural recursion:
\begin{itemize}[nosep]
  \item \textbf{Primitives:} Fixed-width encodings. \texttt{Int} uses
        variable-length encoding (LEB128). \texttt{Float} uses IEEE 754.
        \texttt{String} uses length-prefixed UTF-8. \texttt{Bool} uses
        one byte.
  \item \textbf{Records (products):} Fields are encoded in declaration
        order with field tags (small integers). Each field is prefixed
        with its tag and length:
        $\mathsf{enc}(\{l_1 = v_1, \ldots, l_n = v_n\}) =
        \mathsf{tag}(l_1) \cdot \mathsf{len}(v_1) \cdot \mathsf{enc}(v_1)
        \cdots \mathsf{tag}(l_n) \cdot \mathsf{len}(v_n) \cdot
        \mathsf{enc}(v_n)$.
  \item \textbf{Variants (sums):} A variant tag (small integer)
        followed by the encoding of the payload:
        $\mathsf{enc}(C_i(v)) = \mathsf{tag}(C_i) \cdot \mathsf{enc}(v)$.
  \item \textbf{Containers:} Length prefix followed by element encodings:
        $\mathsf{enc}([v_1, \ldots, v_n]) = n \cdot \mathsf{enc}(v_1)
        \cdots \mathsf{enc}(v_n)$.
\end{itemize}
\end{definition}

\subsection{Schema Evolution}

One of the most challenging aspects of distributed systems is managing
protocol evolution during rolling deployments. At any given time, different
nodes may be running different versions of the code with different type
definitions. \japl{} manages this through type compatibility rules derived
from \Cref{thm:evolution}.

The following changes are backward compatible:

\begin{table}[H]
\centering
\caption{Schema Evolution Rules}
\label{tab:evolution}
\begin{tabularx}{\textwidth}{lXX}
\toprule
\textbf{Change} & \textbf{Rule} & \textbf{Example} \\
\midrule
Add optional field & New field must have \texttt{Option} type
  (defaults to \texttt{None}) &
  \texttt{\{id, name\}} $\to$ \texttt{\{id, name, email: Option<String>\}} \\
\addlinespace
Add variant & New variant added to sum type; old code ignores it &
  \texttt{Active | Inactive} $\to$ \texttt{Active | Inactive | Suspended} \\
\addlinespace
Widen field type & Replace type with strictly more general type &
  \texttt{Int} $\to$ \texttt{Float} (numeric widening) \\
\addlinespace
Rename field & Old tag is preserved; new name is alias &
  \texttt{name} $\to$ \texttt{display\_name} with
  \texttt{@alias("name")} \\
\bottomrule
\end{tabularx}
\end{table}

The following changes are \emph{not} backward compatible and are flagged
by the compiler when it detects version mismatches:

\begin{itemize}[nosep]
  \item Removing a required field.
  \item Removing a variant that may still exist in flight.
  \item Narrowing a field type.
  \item Changing a field tag number.
\end{itemize}

\subsection{Versioned Types}

For cases where incompatible changes are necessary, \japl{} supports
explicit versioning:

\begin{lstlisting}
type Order @version(2) = {
  id: OrderId,
  items: List<Item>,
  total: Money,
  status: OrderStatus,
  -- New in v2
  metadata: Option<Map<String, String>>
}

-- Migration function for v1 -> v2
migrate Order from 1 to 2 = fn(old) ->
  { old | metadata = None }
\end{lstlisting}

The runtime automatically applies migration functions when it receives a
message with an older version tag. This happens transparently, with no
programmer intervention at the receive site.

\subsection{Comparison with External Schema Languages}

\begin{table}[H]
\centering
\caption{Comparison of Serialization Approaches}
\label{tab:serialization}
\begin{tabularx}{\textwidth}{lcccc}
\toprule
\textbf{Property} & \textbf{JAPL} & \textbf{Protobuf} & \textbf{Avro} & \textbf{Cap'n Proto} \\
\midrule
Schema language       & Types & Separate & Separate & Separate \\
Code generation       & No    & Yes      & Yes      & Yes \\
Schema evolution      & Type rules & Field numbers & Schema registry & Ordinals \\
Zero-copy             & Partial & No & No & Yes \\
Type safety           & Full  & Partial  & None     & Partial \\
In-language           & Yes   & No       & No       & No \\
\bottomrule
\end{tabularx}
\end{table}


% ===========================================================================
% 6. LOCATION TRANSPARENCY
% ===========================================================================
\section{Location Transparency and Awareness}
\label{sec:location-transparency}

\subsection{The Transparency Spectrum}

Location transparency is not binary; it exists on a spectrum. At one
extreme, every operation looks identical regardless of location (full
transparency). At the other, every remote operation requires explicit
handling of latency, failure, and serialization (full awareness). \japl{}
occupies a principled position on this spectrum:

\begin{itemize}[nosep]
  \item \textbf{Transparent:} Process send and receive syntax.
  \item \textbf{Aware:} The \texttt{Net} effect, timeouts, partition
        handling, and \texttt{NodeDown} messages.
\end{itemize}

\subsection{Spawn Remote}

The \texttt{Process.spawn\_on} primitive creates a process on a specified
remote node. The returned PID is indistinguishable from a local PID in
subsequent operations:

\begin{lstlisting}
fn deploy_worker(target: Node, config: WorkerConfig)
    -> Pid[WorkerMsg] with Process, Net =
  let pid = Process.spawn_on(target, fn ->
    worker_loop(WorkerState.init(config))
  )
  Process.monitor(pid)
  pid
\end{lstlisting}

The runtime handles the mechanics of remote spawning through a careful
distinction between \emph{code} and \emph{data}. The function passed to
\texttt{spawn\_on} must be a top-level function or a closure over
serializable values. The \emph{code itself is not transferred}: instead,
the compiler emits a module-qualified function name (e.g.,
\texttt{MyApp.Workers.worker\_loop}) that both nodes resolve to the same
compiled code (ensured by the cluster's version management). What is
transferred is the closure's \emph{captured environment}---the values of
all free variables in the closure body. The compiler statically verifies
that every captured value satisfies $\Ser(\tau)$; if any captured variable
has a non-serializable type (e.g., a function type or a file handle), the
program is rejected at compile time. This resolves the apparent tension
with $\neg\Ser(\tau \to \sigma)$: function \emph{values} (closures with
opaque code pointers) cannot be sent as messages, but \texttt{spawn\_on}
transmits a function \emph{reference} (a name) plus serializable captured
state, which the remote node reassembles into a runnable closure.

The returned PID encodes the remote node's address for future communication.

\subsection{Transparent Message Routing}

When a process sends a message to a PID, the runtime checks whether the
PID is local or remote:

\begin{enumerate}[nosep]
  \item \textbf{Local:} The message is placed directly in the target
        process's mailbox. Cost: O(1) pointer enqueue.
  \item \textbf{Remote, connected:} The message is serialized,
        transmitted over the existing TCP connection, deserialized, and
        placed in the remote mailbox. Cost: serialization +
        network latency + deserialization.
  \item \textbf{Remote, disconnected:} The message is buffered (up to a
        configurable limit). If the connection is re-established within
        the timeout, buffered messages are sent. Otherwise, a
        \texttt{NodeDown} message is delivered to monitoring processes.
\end{enumerate}

\subsection{When Transparency Helps}

Location transparency is valuable when the programmer wants to write
code that works identically regardless of deployment topology:

\begin{lstlisting}
-- This function works whether the workers are local or remote
fn distribute_work(workers: List<Pid[WorkerMsg]>, tasks: List<Task>) -> Unit =
  List.zip(workers, tasks)
  |> List.each(fn (pid, task) -> Process.send(pid, DoWork(task)))
\end{lstlisting}

This enables testing distributed code locally (all processes on one node)
and deploying it across multiple nodes without code changes.

\subsection{When Awareness Is Needed}

Some situations require explicit awareness of distribution:

\begin{lstlisting}
-- Explicit timeout for remote operations
fn query_remote_service(service: Pid[QueryMsg], query: Query)
    -> Result<Response, DistError> with Process, Net =
  let ref = Process.monitor(service)
  Process.send(service, DoQuery(query, self()))
  Process.receive_with_timeout(5000) {
  | QueryResult(result) ->
      Process.demonitor(ref)
      Ok(result)
  | ProcessDown(^ref, reason) ->
      Err(ServiceCrashed(reason))
  | Timeout ->
      Process.demonitor(ref)
      Err(ServiceTimeout)
  }
\end{lstlisting}

The \texttt{Net} effect in the function signature signals to callers that
this function involves network communication and may fail in distributed-
specific ways.


% ===========================================================================
% 7. PARTITION TOLERANCE
% ===========================================================================
\section{Partition Tolerance}
\label{sec:partition}

\subsection{The CAP Context}

The CAP theorem~\citep{brewer2000towards,gilbert2002brewer} establishes
that a distributed system cannot simultaneously provide consistency,
availability, and partition tolerance. Since network partitions are
inevitable in practice~\citep{bailis2014network}, distributed systems must
choose between consistency and availability during partitions.

\japl{} does not make this choice for the programmer. Instead, it provides
the primitives to implement both CP and AP strategies, and uses the type
system to make the choice explicit.

\subsection{Process Model and Partitions}

\japl{}'s process model naturally handles network partitions because
processes are already isolated and communicate asynchronously. When a
partition occurs:

\begin{enumerate}[nosep]
  \item Messages to processes on unreachable nodes are buffered.
  \item Monitors trigger \texttt{NodeDown} events for unreachable nodes.
  \item Supervisors can apply partition-specific strategies.
  \item When the partition heals, connections are re-established and
        buffered messages are delivered.
\end{enumerate}

\begin{lstlisting}
-- Handling partitions in a supervisor
fn partition_aware_supervisor() -> Never with Process[SupervisorMsg], Net =
  match Process.receive() with
  | NodeDown(node, reason) ->
      match reason with
      | NetworkPartition ->
          -- Apply split-brain strategy
          let action = SplitBrain.evaluate(cluster_state(), node)
          match action with
          | KeepRunning -> partition_aware_supervisor()
          | StepDown -> Process.exit(PartitionStepDown)
          | FenceNode(n) -> Node.disconnect(n)
      | NodeCrash ->
          -- Restart children that were on the crashed node
          restart_children_on(available_node())
          partition_aware_supervisor()
  | PartitionHealed(node) ->
      -- Reconcile state with the re-connected node
      reconcile_state(node)
      partition_aware_supervisor()
  | msg ->
      handle_normal(msg)
      partition_aware_supervisor()
\end{lstlisting}

\subsection{Split-Brain Strategies}

\japl{} provides built-in split-brain resolution strategies:

\begin{description}[nosep]
  \item[Static quorum:] A partition is viable if it contains a majority
    of configured nodes. Minority partitions shut down.
  \item[Keep oldest:] The partition containing the oldest node survives.
  \item[Keep referee:] A designated referee node decides which partition
    survives.
  \item[Custom:] The programmer provides a function
    $\texttt{fn(ClusterState) -> SplitBrainAction}$.
\end{description}

\begin{lstlisting}
-- Configuring split-brain strategy
let cluster = Cluster.start({
  name = "payment-service",
  nodes = seed_nodes,
  split_brain = SplitBrain.StaticQuorum({ quorum_size = 3 }),
  rejoin = Rejoin.CleanStart,
})
\end{lstlisting}

\subsection{CRDTs for Partition-Tolerant State}

For state that must remain available during partitions, \japl{} provides
CRDT (Conflict-free Replicated Data Type) primitives~\citep{shapiro2011comprehensive}:

\begin{lstlisting}
-- A grow-only counter replicated across nodes
type DistributedCounter = crdt GCounter

fn increment_counter(counter: DistributedCounter) -> DistributedCounter =
  GCounter.increment(counter, Node.self())

-- CRDTs merge automatically when partitions heal
fn on_partition_heal(local: DistributedCounter, remote: DistributedCounter)
    -> DistributedCounter =
  GCounter.merge(local, remote)  -- always converges, no conflicts
\end{lstlisting}


% ===========================================================================
% 8. SERVICE DISCOVERY AND CLUSTERING
% ===========================================================================
\section{Service Discovery and Clustering}
\label{sec:discovery}

\subsection{Node Mesh}

\japl{} nodes form a mesh network with full connectivity (every node
maintains a connection to every other node, up to a configurable maximum).
The mesh is established and maintained through a gossip-based membership
protocol~\citep{vanrenesse1998gossip}.

\begin{lstlisting}
-- Starting a node with mesh configuration
let node = Node.start({
  name = "api-1",
  cookie = Env.get("CLUSTER_COOKIE"),
  listen = "0.0.0.0:9000",
  seeds = ["seed-1.internal:9000", "seed-2.internal:9000"],
  max_connections = 100,
})
\end{lstlisting}

\subsection{Membership Protocol}

The membership protocol ensures that all nodes have a consistent view of
cluster membership:

\begin{enumerate}[nosep]
  \item \textbf{Join:} A new node connects to a seed node and sends a
        join request. The seed gossips the new member to all other nodes.
  \item \textbf{Heartbeat:} Nodes exchange heartbeats at regular intervals.
        The phi accrual failure detector~\citep{hayashibara2004phi}
        determines when a node is unreachable.
  \item \textbf{Leave:} A node announces its intention to leave. Other
        nodes remove it from their membership list gracefully.
  \item \textbf{Down:} If a node fails to respond within the failure
        detection threshold, it is marked as down. Monitors are notified.
\end{enumerate}

\subsection{Service Registry}

The service registry is a distributed key-value store mapping service
names to sets of PIDs. It is replicated across all nodes using a CRDT
(OR-Set):

\begin{lstlisting}
-- Register the current process as a service
Registry.register("order-processor", self())

-- Register with metadata
Registry.register("order-processor", self(), {
  version = "2.1.0",
  region = "us-east-1",
  capacity = 100,
})

-- Look up all instances of a service
let processors: List<{pid: Pid[OrderMsg], meta: ServiceMeta}> =
  Registry.lookup("order-processor")

-- Look up with filter
let local_processors =
  Registry.lookup("order-processor")
  |> List.filter(fn entry -> entry.meta.region == Node.region())

-- Subscribe to changes
Registry.subscribe("order-processor", fn event ->
  match event with
  | ServiceRegistered(pid, meta) -> log("New processor: " ++ show(pid))
  | ServiceDeregistered(pid) -> log("Processor left: " ++ show(pid))
)
\end{lstlisting}

\subsection{Health Checking}

\japl{} provides built-in health checking at two levels:

\begin{description}[nosep]
  \item[Node-level:] The failure detector monitors node connectivity.
    Unreachable nodes trigger \texttt{NodeDown} events.
  \item[Service-level:] Registered services can define health check
    functions that are invoked periodically:
\end{description}

\begin{lstlisting}
Registry.register("database-pool", self(), {
  health_check = fn pid ->
    let reply = Process.call(pid, HealthCheck, 5000)
    match reply with
    | Ok(Healthy) -> HealthStatus.Healthy
    | Ok(Degraded(reason)) -> HealthStatus.Degraded(reason)
    | Err(Timeout) -> HealthStatus.Unhealthy("timeout")
})
\end{lstlisting}

\subsection{Load Balancing}

The service registry integrates with load balancing strategies:

\begin{lstlisting}
-- Round-robin selection from registered services
let pid = Registry.select("order-processor", Strategy.RoundRobin)

-- Least-loaded selection (requires capacity metadata)
let pid = Registry.select("order-processor", Strategy.LeastLoaded)

-- Locality-aware: prefer local node, then same region, then any
let pid = Registry.select("order-processor", Strategy.LocalFirst)
\end{lstlisting}


% ===========================================================================
% 9. COMPARISON
% ===========================================================================
\section{Comparison with Existing Approaches}
\label{sec:comparison}

\begin{table}[H]
\centering
\caption{Distribution Feature Comparison}
\label{tab:comparison}
\small
\begin{tabularx}{\textwidth}{lccccccX}
\toprule
\textbf{Feature} & \rotatebox{70}{\textbf{JAPL}} &
\rotatebox{70}{\textbf{Erlang}} & \rotatebox{70}{\textbf{Akka}} &
\rotatebox{70}{\textbf{Orleans}} & \rotatebox{70}{\textbf{Go}} &
\rotatebox{70}{\textbf{Rust}} & \rotatebox{70}{\textbf{Unison}} \\
\midrule
Native distribution          & \checkmark & \checkmark & Lib & Lib & -- & -- & \checkmark \\
Location transparency        & \checkmark & \checkmark & \checkmark & \checkmark & -- & -- & \checkmark \\
Static typing                & \checkmark & --         & \checkmark & \checkmark & \checkmark & \checkmark & \checkmark \\
Type-derived serialization   & \checkmark & --         & --         & Partial & -- & -- & Hash \\
Schema evolution             & \checkmark & --         & Manual     & Manual  & -- & -- & Hash \\
Built-in supervision         & \checkmark & \checkmark & \checkmark & --      & -- & -- & -- \\
Distributed supervision      & \checkmark & \checkmark & \checkmark & --      & -- & -- & -- \\
Service discovery            & \checkmark & EPMD       & Config     & Silo    & -- & -- & -- \\
Effect-marked distribution   & \checkmark & --         & --         & --      & -- & -- & -- \\
Capability security          & \checkmark & --         & --         & --      & -- & -- & -- \\
CRDTs built-in               & \checkmark & Lib        & Lib        & --      & -- & -- & -- \\
Partition strategies          & \checkmark & Lib        & \checkmark & --      & -- & -- & -- \\
\bottomrule
\end{tabularx}
\end{table}

\subsection{Erlang: The Gold Standard}

Erlang is the closest language to \japl{} in distribution philosophy. Both
treat distribution as native, both use location-transparent PIDs, and both
provide supervision trees. The key differences are:

\begin{enumerate}[nosep]
  \item \textbf{Type safety.} Erlang is dynamically typed. A serialization
        error in Erlang manifests as a runtime crash. In \japl{}, the
        compiler rejects programs that attempt to send non-serializable
        values across nodes.
  \item \textbf{Schema evolution.} Erlang's External Term Format does not
        support schema evolution. Adding a field to a record requires
        coordinating all nodes. \japl{}'s type compatibility rules allow
        independent evolution.
  \item \textbf{Effect tracking.} Erlang does not distinguish local from
        distributed operations in types. \japl{}'s \texttt{Net} effect
        makes this distinction visible.
  \item \textbf{Security.} Erlang's cookie-based authentication provides
        all-or-nothing access. \japl{}'s capability types allow
        fine-grained distributed access control.
\end{enumerate}

\subsection{Akka: Library-Level Actors on the JVM}

Akka demonstrates that actor-based distribution can work as a library, but
also reveals the costs:

\begin{itemize}[nosep]
  \item Serialization must be configured externally (via
        \texttt{application.conf} or annotations).
  \item Cluster membership uses configuration files and seed nodes, not
        language primitives.
  \item The JVM's exception model interacts poorly with actor supervision.
  \item Akka Cluster is a sophisticated system, but its complexity (split
        brain resolvers, cluster sharding, distributed data) reflects the
        difficulty of building distribution atop a language not designed
        for it.
\end{itemize}

\subsection{Orleans: Virtual Actors}

Orleans' virtual actor model is elegant: actors always exist conceptually
and are activated on demand. This eliminates lifecycle management but
introduces other complexities:

\begin{itemize}[nosep]
  \item Grain interfaces require code generation (separate from the
        language's type system).
  \item Persistence requires explicit configuration with external storage.
  \item The virtual actor model assumes that actors are stateless between
        activations, which does not fit all patterns.
\end{itemize}

\japl{} could implement a virtual-actor pattern as a library atop its
native process model, providing the best of both worlds.

\subsection{Go: No Built-In Distribution}

Go provides goroutines and channels for local concurrency but offers nothing
for distribution. Distributed Go programs use gRPC (with Protocol Buffers),
HTTP APIs, or messaging systems like NATS. This means every distributed Go
program must:

\begin{enumerate}[nosep]
  \item Define Protocol Buffer schemas separately from Go types.
  \item Generate Go code from the schemas.
  \item Handle serialization errors at runtime.
  \item Implement service discovery using external systems (Consul, etcd).
  \item Implement health checking and failure detection manually.
\end{enumerate}

\subsection{Rust: No Built-In Distribution}

Rust's ownership model provides excellent local correctness guarantees but
does not extend to distribution. Distributed Rust programs face the same
challenges as Go, with the additional complexity of lifetime management
across network boundaries. Libraries like Actix provide actor-like
abstractions, but serialization (via Serde) is a separate concern.

\subsection{Unison: Radical Content-Addressing}

Unison's content-addressed approach is the most radical departure from
traditional distribution. By identifying code by hash, Unison eliminates
versioning and serialization problems entirely---any value can be sent
to any node, and if the code is not present, it can be transmitted.

However, Unison's approach requires abandoning file-based source code,
traditional version control, and familiar development workflows. \japl{}
takes a more conservative approach that integrates with existing tools
while still providing strong distribution guarantees.


% ===========================================================================
% 10. IMPLEMENTATION
% ===========================================================================
\section{Implementation Architecture}
\label{sec:implementation}

\subsection{Network Layer}

The network layer is organized in four tiers:

\begin{figure}[H]
\centering
\begin{tikzpicture}[
  layer/.style={draw, rounded corners, minimum width=12cm, minimum height=1.2cm,
                fill=blue!5, font=\small},
  arr/.style={-Stealth, thick},
]
  \node[layer] (app)  at (0, 4.5) {Application Layer: Process.send / Process.spawn\_on / Registry};
  \node[layer] (dist) at (0, 3.0) {Distribution Layer: PID routing, serialization, buffering};
  \node[layer] (conn) at (0, 1.5) {Connection Layer: connection pool, multiplexing, heartbeat};
  \node[layer] (tcp)  at (0, 0)   {Transport Layer: TCP/TLS, QUIC (optional)};

  \draw[arr] (app.south) -- (dist.north);
  \draw[arr] (dist.south) -- (conn.north);
  \draw[arr] (conn.south) -- (tcp.north);
\end{tikzpicture}
\caption{Network layer architecture.}
\label{fig:network-layers}
\end{figure}

\subsection{Connection Pooling}

Each node maintains a connection pool to every connected peer. Connections
are multiplexed: a single TCP connection carries messages for all processes
communicating between two nodes. This avoids the overhead of per-process
connections while allowing fair scheduling across message streams.

The connection pool supports:
\begin{itemize}[nosep]
  \item \textbf{Automatic reconnection} with exponential backoff.
  \item \textbf{Connection migration} from TCP to QUIC when both nodes support it.
  \item \textbf{Flow control} per message stream to prevent one chatty
        process from starving others.
  \item \textbf{Priority queues} for supervision messages (which must be
        delivered with minimal latency).
\end{itemize}

\subsection{Message Encoding}

Messages are encoded using the type-derived wire format described in
\Cref{sec:wire-protocols}. The encoding pipeline is:

\begin{enumerate}[nosep]
  \item \textbf{Serialization:} The value is serialized to bytes using the
        type-derived encoder.
  \item \textbf{Framing:} The serialized bytes are wrapped in a frame with
        a header containing: message type tag, version, source PID,
        destination PID, and length.
  \item \textbf{Compression:} For messages above a configurable threshold
        (default: 1024 bytes), LZ4 compression is applied.
  \item \textbf{Encryption:} If the connection uses TLS, encryption is
        handled at the transport layer. For additional message-level
        encryption, the capability system can require encrypted channels.
\end{enumerate}

The frame format is:

\begin{verbatim}
+--------+--------+--------+--------+--------+--------+
| Magic  | Flags  | MsgType| Version| Length           |
| (2B)   | (1B)   | (4B)   | (2B)   | (4B)            |
+--------+--------+--------+--------+--------+--------+
| Source PID (16B)          | Dest PID (16B)           |
+--------+--------+--------+--------+--------+--------+
| Capability Token (0 or 32B, if Flags & 0x10)         |
+--------+--------+--------+--------+--------+--------+
| Payload (variable length)                             |
+--------+--------+--------+--------+--------+--------+
\end{verbatim}

\subsection{Node Discovery Protocol}

Node discovery uses a two-phase protocol:

\begin{description}[nosep]
  \item[Phase 1: Seed contact.] The joining node connects to one or more
    seed nodes from its configuration. It sends a \texttt{JoinRequest}
    containing its name, capabilities, and listening address.
  \item[Phase 2: Gossip dissemination.] The seed node adds the new member
    to its membership list and gossips the change to all other nodes.
    Each node then establishes a direct connection to the new member.
\end{description}

The gossip protocol uses a push-pull model: each gossip round, a node
selects a random peer, sends its membership state, and receives the
peer's state. States are merged using a vector clock for convergence.

\subsection{Failure Detection: Phi Accrual}

\japl{} uses the phi accrual failure detector~\citep{hayashibara2004phi}
to determine node unreachability. Unlike binary failure detectors (alive
or dead), the phi accrual detector outputs a continuous \emph{suspicion
level} $\phi$ that represents the likelihood that a node has failed.

The suspicion level is computed from the distribution of inter-heartbeat
intervals:

\[
  \phi(t) = -\log_{10}\left(1 - F(t - t_{\text{last}})\right)
\]

where $F$ is the cumulative distribution function of the normal
distribution fitted to historical heartbeat intervals, $t$ is the current
time, and $t_{\text{last}}$ is the time of the last received heartbeat.

Advantages of the phi accrual detector:
\begin{itemize}[nosep]
  \item Adapts to network conditions (high-latency links produce higher
        thresholds automatically).
  \item Produces fewer false positives than fixed-timeout detectors.
  \item The threshold $\phi_{\text{max}}$ is configurable per application
        (lower for latency-sensitive, higher for stability-sensitive).
\end{itemize}


% ===========================================================================
% 11. SECURITY
% ===========================================================================
\section{Distributed Security}
\label{sec:security}

\subsection{Capability-Based Security Model}

\japl{}'s distribution model uses capabilities---unforgeable tokens that
represent permission to perform specific operations---as the foundation for
distributed security. This approach is inspired by the E
language~\citep{miller2006robust} and the object-capability
model~\citep{dennis1966programming,miller2003capability}.

A PID and a capability are distinct concepts. A \texttt{Pid[Msg]} is a
location-transparent process address that enables message delivery; possessing
a PID allows sending messages of type \texttt{Msg} to the identified process,
but it does \emph{not} grant permission to spawn processes on that node,
access the service registry, or perform other privileged operations. For those,
a separate \emph{capability token} must be obtained. Capability tokens are
cryptographically signed, unforgeable values that encode a specific permission
scope. They are first-class values that can be passed in messages (they
implement $\Ser$), attenuated, but never amplified.

\begin{lstlisting}
-- A capability to spawn processes on a specific node.
-- This is a separate token, distinct from any PID.
type SpawnCapability = capability {
  node: Node,
  max_processes: Int,
  allowed_types: Set<TypeId>,
  token: CapabilityToken,  -- cryptographic proof of authority
}

-- Using the capability
fn deploy_service(cap: SpawnCapability, service: fn() -> Never)
    -> Result<Pid[ServiceMsg], CapError> with Process, Net =
  if SpawnCapability.check(cap, type_of(service)) then
    Ok(Process.spawn_on(cap.node, service))
  else
    Err(NotAllowed)
\end{lstlisting}

\noindent The \texttt{SpawnCapability.check} call is a \emph{compile-time}
verification of type compatibility (ensuring the service type is in
\texttt{allowed\_types}) combined with a \emph{runtime} verification of
the cryptographic token's validity and the process count budget. The
compile-time component is a static check inserted by the compiler; it
rejects programs where the service type is provably outside the capability's
\texttt{allowed\_types}. The runtime component validates the token signature
and decrements the process counter, which cannot be checked statically
because it depends on dynamic cluster state. This two-phase design ensures
that no runtime latency is added for checks that can be resolved
statically, while dynamic invariants are still enforced.

\subsection{Node Authentication}

When two \japl{} nodes establish a connection, they perform mutual
authentication:

\begin{enumerate}[nosep]
  \item \textbf{TLS handshake:} All inter-node connections use TLS 1.3
        with mutual certificate verification.
  \item \textbf{Cluster cookie:} After TLS, nodes exchange a shared
        secret (the cluster cookie) as an additional authentication layer.
  \item \textbf{Capability exchange:} Nodes exchange their capability
        sets, establishing what operations each can perform on the other.
\end{enumerate}

\begin{lstlisting}
let node = Node.start({
  name = "api-1",
  cookie = Env.get("CLUSTER_COOKIE"),
  listen = "0.0.0.0:9000",
  tls = TlsConfig.from_files(
    cert = "/etc/japl/node.crt",
    key = "/etc/japl/node.key",
    ca = "/etc/japl/ca.crt",
  ),
  capabilities = NodeCapabilities.from([
    Cap.SpawnProcesses({ max = 1000 }),
    Cap.SendMessages({ rate_limit = 10000 }),
    Cap.AccessRegistry({ read = True, write = False }),
  ]),
})
\end{lstlisting}

\subsection{Encrypted Channels}

For sensitive messages, \japl{} supports end-to-end encrypted channels
that are encrypted beyond the TLS transport layer:

\begin{lstlisting}
-- Create an encrypted channel between two processes
let (send_chan, recv_chan) = Channel.encrypted(
  algorithm = Aes256Gcm,
  key = shared_key,
)

-- Messages on this channel are encrypted before serialization
Channel.send(send_chan, SensitiveData(patient_record))
\end{lstlisting}

\subsection{Distributed Access Control}

Capabilities can be attenuated (reduced in power) but never amplified.
This follows the principle of least authority~\citep{miller2006robust}:

\begin{lstlisting}
-- A full capability
let full_cap = RegistryCapability.full()

-- Attenuate to read-only
let read_cap = RegistryCapability.attenuate(full_cap, {
  read = True,
  write = False,
  subscribe = True,
})

-- Attenuate further to a specific service
let limited_cap = RegistryCapability.attenuate(read_cap, {
  service_filter = fn name -> name == "payment-service",
})

-- Send the limited capability to a remote process
Process.send(remote_pid, GrantAccess(limited_cap))
\end{lstlisting}


% ===========================================================================
% 12. CASE STUDY
% ===========================================================================
\section{Case Study: Distributed Task Scheduler}
\label{sec:case-study}

We demonstrate \japl{}'s distribution model through a complete distributed
task scheduler that manages work across multiple nodes, including process
migration when nodes become overloaded.

\subsection{System Architecture}

The task scheduler consists of three components:
\begin{enumerate}[nosep]
  \item A \textbf{coordinator} process that receives task submissions and
        assigns them to worker nodes.
  \item \textbf{Worker pools} on each node that execute tasks.
  \item A \textbf{monitor} that tracks node load and triggers process
        migration.
\end{enumerate}

\begin{figure}[H]
\centering
\begin{tikzpicture}[
  node_box/.style={draw, rounded corners, minimum width=3cm, minimum height=2cm,
                   fill=blue!5},
  proc/.style={draw, circle, minimum size=0.6cm, fill=green!20, font=\tiny},
  arr/.style={-Stealth, thick},
]
  % Coordinator node
  \node[node_box, label=above:{Node: coordinator}] (coord) at (0, 0) {};
  \node[proc] (c1) at (-0.5, 0.3) {C};
  \node[proc] (c2) at (0.5, 0.3) {M};
  \node[font=\tiny] at (0, -0.5) {Coordinator + Monitor};

  % Worker node 1
  \node[node_box, label=above:{Node: worker-1}] (w1) at (-4.5, -3.5) {};
  \node[proc] (w1a) at (-5.2, -3.2) {W};
  \node[proc] (w1b) at (-4.5, -3.2) {W};
  \node[proc] (w1c) at (-3.8, -3.2) {W};
  \node[font=\tiny] at (-4.5, -4.0) {Worker Pool};

  % Worker node 2
  \node[node_box, label=above:{Node: worker-2}] (w2) at (0, -3.5) {};
  \node[proc] (w2a) at (-0.7, -3.2) {W};
  \node[proc] (w2b) at (0, -3.2) {W};
  \node[proc] (w2c) at (0.7, -3.2) {W};
  \node[font=\tiny] at (0, -4.0) {Worker Pool};

  % Worker node 3
  \node[node_box, label=above:{Node: worker-3}] (w3) at (4.5, -3.5) {};
  \node[proc] (w3a) at (3.8, -3.2) {W};
  \node[proc] (w3b) at (4.5, -3.2) {W};
  \node[proc] (w3c) at (5.2, -3.2) {W};
  \node[font=\tiny] at (4.5, -4.0) {Worker Pool};

  % Arrows
  \draw[arr] (coord.south west) -- (w1.north);
  \draw[arr] (coord.south) -- (w2.north);
  \draw[arr] (coord.south east) -- (w3.north);

  % Migration arrow
  \draw[arr, dashed, red, thick] (w1.east) -- node[above, font=\tiny, red] {migrate} (w2.west);
\end{tikzpicture}
\caption{Distributed task scheduler architecture with process migration.}
\label{fig:scheduler}
\end{figure}

\subsection{Message Types}

\begin{lstlisting}
type TaskId = String

type Task deriving(Serialize) = {
  id: TaskId,
  payload: Bytes,
  priority: Priority,
  deadline: Option<Timestamp>,
}

type Priority = High | Medium | Low
  deriving(Serialize, Ord)

type CoordinatorMsg =
  | SubmitTask(Task, Reply<Result<TaskId, SchedulerError>>)
  | TaskCompleted(TaskId, TaskResult)
  | WorkerOverloaded(Node, Float)
  | NodeJoined(Node)
  | NodeLeft(Node)

type WorkerMsg =
  | ExecuteTask(Task)
  | MigrateTask(Task, Node)
  | Shutdown

type MonitorMsg =
  | LoadReport(Node, Float)
  | CheckLoads
\end{lstlisting}

\subsection{Coordinator Implementation}

\begin{lstlisting}
type SchedulerState = {
  workers: Map<Node, List<Pid<WorkerMsg>>>,
  pending: Map<TaskId, Task>,
  load: Map<Node, Float>,
  round_robin_idx: Int,
}

fn coordinator(state: SchedulerState) -> Never
    with Process[CoordinatorMsg], Net =
  match Process.receive() with
  | SubmitTask(task, reply) ->
      -- Select least-loaded node
      let target = select_node(state)
      match target with
      | Some(node) ->
          let worker = select_worker(state.workers, node)
          Process.send(worker, ExecuteTask(task))
          Reply.send(reply, Ok(task.id))
          let new_pending = Map.insert(state.pending, task.id, task)
          coordinator({ state | pending = new_pending })
      | None ->
          Reply.send(reply, Err(NoWorkersAvailable))
          coordinator(state)

  | TaskCompleted(task_id, result) ->
      let new_pending = Map.delete(state.pending, task_id)
      log("Task " ++ task_id ++ " completed: " ++ show(result))
      coordinator({ state | pending = new_pending })

  | WorkerOverloaded(node, load) ->
      -- Trigger migration from overloaded node
      let tasks_to_migrate = select_tasks_for_migration(state, node)
      let target = least_loaded_node(state, node)
      match target with
      | Some(target_node) ->
          List.each(tasks_to_migrate, fn task ->
            let worker = select_worker(state.workers, node)
            Process.send(worker, MigrateTask(task, target_node))
          )
      | None -> ()
      coordinator({ state | load = Map.insert(state.load, node, load) })

  | NodeJoined(node) ->
      -- Spawn workers on the new node
      let pids = List.range(1, 4)
        |> List.map(fn _ ->
             Process.spawn_on(node, fn -> worker_loop(WorkerState.init()))
           )
      let new_workers = Map.insert(state.workers, node, pids)
      coordinator({ state | workers = new_workers })

  | NodeLeft(node) ->
      -- Reassign tasks from the departed node
      let orphaned = get_tasks_on_node(state, node)
      List.each(orphaned, fn task ->
        Process.send(self(), SubmitTask(task, Reply.ignore()))
      )
      let new_workers = Map.delete(state.workers, node)
      coordinator({ state | workers = new_workers })
\end{lstlisting}

\subsection{Worker Implementation}

\begin{lstlisting}
type WorkerState = {
  current_task: Option<Task>,
  tasks_completed: Int,
}

fn worker_loop(state: WorkerState) -> Never with Process[WorkerMsg] =
  match Process.receive() with
  | ExecuteTask(task) ->
      let result = execute(task)
      -- Report completion to coordinator
      Process.send(coordinator_pid(), TaskCompleted(task.id, result))
      worker_loop({ state |
        current_task = None,
        tasks_completed = state.tasks_completed + 1
      })

  | MigrateTask(task, target_node) ->
      -- Spawn a new worker on the target node with the task
      let new_worker = Process.spawn_on(target_node, fn ->
        let result = execute(task)
        Process.send(coordinator_pid(), TaskCompleted(task.id, result))
        worker_loop(WorkerState.init())
      )
      Process.monitor(new_worker)
      worker_loop({ state | current_task = None })

  | Shutdown ->
      log("Worker shutting down after " ++
          show(state.tasks_completed) ++ " tasks")
      Process.exit(Normal)
\end{lstlisting}

\paragraph{Mailbox forwarding during migration.}
When a process ``migrates'' by spawning a replacement on a target node, the
original PID remains valid. The \japl{} runtime installs a \emph{forwarding
entry} in the local node's routing table: any message arriving for the old
PID is transparently re-routed to the new PID on the target node. This
forwarding is invisible to senders, preserving location transparency. The
forwarding entry persists until (a)~all monitors of the old PID have
received a \texttt{ProcessMigrated} notification and updated their
references, or (b)~a configurable forwarding TTL expires, after which
messages to the old PID result in a \texttt{ProcessNotFound} error
delivered to the sender's monitor. In the task scheduler above, the
coordinator continues to hold the old worker PID until it receives the
\texttt{TaskCompleted} message from the new worker; at that point it can
update its routing table. For long-lived migrations, the runtime's
forwarding mechanism ensures no messages are lost during the transition
window.

\subsection{Load Monitor}

\begin{lstlisting}
fn load_monitor(state: MonitorState) -> Never
    with Process[MonitorMsg], Net =
  match Process.receive_with_timeout(5000) with
  | LoadReport(node, load) ->
      let new_state = { state |
        loads = Map.insert(state.loads, node, load)
      }
      -- Check if any node is overloaded
      if load > state.threshold then
        Process.send(state.coordinator, WorkerOverloaded(node, load))
      load_monitor(new_state)

  | CheckLoads ->
      -- Periodic check: request load from all nodes
      Map.keys(state.loads)
      |> List.each(fn node ->
           let load = Node.get_load(node)
           Process.send(self(), LoadReport(node, load))
         )
      load_monitor(state)

  | Timeout ->
      -- Timeout triggers periodic check
      Process.send(self(), CheckLoads)
      load_monitor(state)
\end{lstlisting}

\subsection{Supervisor Tree}

\begin{lstlisting}
supervisor TaskScheduler {
  strategy: OneForOne
  max_restarts: 10
  max_seconds: 60

  child {
    id: "coordinator"
    start: fn -> coordinator(SchedulerState.init())
    restart: Permanent
  }

  child {
    id: "load_monitor"
    start: fn -> load_monitor(MonitorState.init())
    restart: Permanent
  }

  child {
    id: "worker_supervisor"
    start: fn -> worker_supervisor()
    restart: Permanent
  }
}

supervisor WorkerSupervisor {
  strategy: OneForOne

  -- Dynamic children: workers are added as nodes join
  dynamic: True
}
\end{lstlisting}

\subsection{Starting the Distributed System}

\begin{lstlisting}
fn main() -> Unit with Io, Net, Process =
  -- Start the node
  let node = Node.start({
    name = "scheduler",
    cookie = Env.get("CLUSTER_COOKIE"),
    listen = "0.0.0.0:9000",
    seeds = Env.get("SEED_NODES") |> String.split(","),
  })

  -- Subscribe to cluster events
  Cluster.subscribe(fn event ->
    match event with
    | MemberJoined(node) ->
        Process.send(coordinator_pid(), NodeJoined(node))
    | MemberLeft(node) ->
        Process.send(coordinator_pid(), NodeLeft(node))
    | _ -> ()
  )

  -- Start the supervisor tree
  let sup = Supervisor.start(TaskScheduler)

  -- The system is now running. Block until shutdown.
  Process.receive_matching(fn msg ->
    match msg with | ShutdownSystem -> True | _ -> False
  )
  Supervisor.stop(sup)
\end{lstlisting}

\subsection{Analysis}

This case study demonstrates several advantages of \japl{}'s native
distribution:

\begin{enumerate}[nosep]
  \item \textbf{No serialization code.} All message types are
        automatically serializable via \texttt{deriving(Serialize)}.
        No Protocol Buffer schemas, no code generation.
  \item \textbf{Uniform local/remote syntax.} The coordinator sends
        messages to workers using the same \texttt{Process.send} regardless
        of whether the worker is local or remote.
  \item \textbf{Process migration.} Migrating a task to another node is
        simply spawning a new process on the target node---no special
        migration framework needed.
  \item \textbf{Fault tolerance.} Supervision trees handle worker crashes,
        node departures, and coordinator failures.
  \item \textbf{Service discovery.} Cluster membership events are delivered
        as messages, integrating naturally with the process model.
\end{enumerate}

A comparable system in Go or Rust would require approximately 2--3x more
code, including Protocol Buffer definitions, gRPC service definitions,
connection management, health checking, and explicit serialization/
deserialization at every network boundary.


% ===========================================================================
% 13. DISCUSSION
% ===========================================================================
\section{Discussion}
\label{sec:discussion}

\subsection{Limitations of Location Transparency}

While \japl{}'s location transparency reduces syntactic burden, it does
not eliminate the fundamental differences between local and remote
communication~\citep{waldo1996note}. The \texttt{Net} effect mitigates
this by making network operations visible in types, but the programmer
must still reason about latency, partial failure, and ordering.

A potential future direction is a \emph{latency-aware type system} where
the type of a remote operation encodes expected latency bounds, allowing
the compiler to warn about operations that may be too slow.

\subsection{Serialization Overhead}

Type-derived serialization is convenient but may not match the performance
of hand-optimized binary protocols. For extremely high-throughput scenarios
(millions of messages per second between two nodes), \japl{} allows
dropping down to custom serialization:

\begin{lstlisting}
-- Custom serializer for a performance-critical type
impl Serialize[MarketTick] =
  fn serialize(tick) =
    -- Hand-optimized packed encoding
    Bytes.builder()
    |> Bytes.write_u64(tick.timestamp)
    |> Bytes.write_u32(tick.symbol_id)
    |> Bytes.write_f64(tick.price)
    |> Bytes.build()

  fn deserialize(data) =
    let timestamp = Bytes.read_u64(data, 0)
    let symbol_id = Bytes.read_u32(data, 8)
    let price = Bytes.read_f64(data, 12)
    Ok({ timestamp, symbol_id, price })
\end{lstlisting}

\subsection{Comparison with Microservice Architectures}

\japl{}'s distribution model can be seen as an alternative to
microservice architectures. Instead of deploying many independent
services that communicate over HTTP/gRPC, a \japl{} application deploys
as a cluster of nodes running a single program. The advantages are:

\begin{itemize}[nosep]
  \item \textbf{Type safety across boundaries.} Inter-service communication
        in microservices is typically untyped at the language level.
        \japl{}'s type system spans node boundaries.
  \item \textbf{No deployment ceremony.} Adding a new node is joining
        the cluster, not deploying a new service.
  \item \textbf{Shared supervision.} A single supervision tree spans the
        entire cluster, providing unified fault tolerance.
\end{itemize}

The disadvantage is reduced independence: all nodes must run compatible
versions of the code. \japl{}'s schema evolution rules mitigate this for
rolling deployments.

\subsection{Scaling Considerations}

The full mesh topology works well for clusters up to approximately 100--200
nodes. Beyond this, \japl{} supports hierarchical clustering where nodes
form groups with full mesh within groups and delegate connections between
groups. This is similar to Erlang's hidden nodes or Akka's cluster
multi-DC support.

\subsection{Formal Verification Opportunities}

The combination of typed processes, effect tracking, and type-derived
serialization opens opportunities for formal verification of distributed
protocols. A \japl{} program can be translated to a process algebra
model and model-checked for properties like deadlock freedom, liveness,
and eventual consistency. We leave this as future work.


% ===========================================================================
% 14. CONCLUSION
% ===========================================================================
\section{Conclusion}
\label{sec:conclusion}

We have presented \japl{}'s approach to making distribution a native
language concern. By integrating location-transparent process identifiers,
type-derived wire protocols, effect-tracked network operations, distributed
supervision, and built-in service discovery into the language, \japl{}
eliminates the ``distribution tax'' that plagues most programming languages.

The key insight is that types are not merely local correctness artifacts;
they are distributed protocol specifications. An algebraic data type
definition simultaneously specifies the in-memory representation, the wire
format, and the schema evolution rules. This unification eliminates the
impedance mismatch between the language and the network.

Our formal framework, based on a location-aware $\pi$-calculus with
type-safe serialization, provides rigorous foundations for reasoning about
distributed \japl{} programs. The categorical semantics using fibered
categories over network topology captures the essential structure of
distributed computation: local fibers of processes over a global base of
network topology.

The case study demonstrates that \japl{}'s native distribution support
reduces distributed systems code by 40--60\% compared to library-based
approaches while providing stronger correctness guarantees. The combination
of static type safety, location transparency, and practical distributed
systems primitives is, to our knowledge, unique among production-oriented
programming languages.

Distribution is not a library concern. It is a language concern. \japl{}
treats it as such.


% ===========================================================================
% REFERENCES
% ===========================================================================
\bibliographystyle{plainnat}

\begin{thebibliography}{40}

\bibitem[Akka(2009)]{akka2009}
Lightbend.
\newblock Akka: Build powerful reactive, concurrent, and distributed
  applications more easily.
\newblock \url{https://akka.io}, 2009.

\bibitem[Alvaro et~al.(2011)]{alvaro2011consistency}
P.~Alvaro, N.~Conway, J.~M. Hellerstein, and W.~R. Marczak.
\newblock Consistency analysis in {Bloom}: A {CALM} and collected approach.
\newblock In \emph{Proc.\ CIDR}, 2011.

\bibitem[Armstrong(2003)]{armstrong2003making}
J.~Armstrong.
\newblock \emph{Making reliable distributed systems in the presence of
  software errors}.
\newblock PhD thesis, Royal Institute of Technology, Stockholm, 2003.

\bibitem[Armstrong(2007)]{armstrong2007programming}
J.~Armstrong.
\newblock \emph{Programming Erlang: Software for a Concurrent World}.
\newblock Pragmatic Bookshelf, 2007.

\bibitem[Avro(2009)]{avro2009}
Apache Foundation.
\newblock Apache {Avro}: A data serialization system.
\newblock \url{https://avro.apache.org}, 2009.

\bibitem[Bailis and Kingsbury(2014)]{bailis2014network}
P.~Bailis and K.~Kingsbury.
\newblock The network is reliable.
\newblock \emph{ACM Queue}, 12(7):15, 2014.

\bibitem[Bernstein et~al.(2014)]{bernstein2014orleans}
P.~A. Bernstein, S.~Bykov, A.~Geller, G.~Kliot, and J.~Thelin.
\newblock Orleans: Distributed virtual actors for programmability and
  scalability.
\newblock Technical Report MSR-TR-2014-41, Microsoft Research, 2014.

\bibitem[Brewer(2000)]{brewer2000towards}
E.~A. Brewer.
\newblock Towards robust distributed systems.
\newblock In \emph{Proc.\ ACM PODC}, 2000.
\newblock Invited talk.

\bibitem[Bykov et~al.(2011)]{bykov2011orleans}
S.~Bykov, A.~Geller, G.~Kliot, J.~R. Larus, R.~Pandya, and J.~Thelin.
\newblock Orleans: Cloud computing for everyone.
\newblock In \emph{Proc.\ ACM SoCC}, pages 16:1--16:14, 2011.

\bibitem[Cap'n Proto(2013)]{capnproto2013}
K.~Varda.
\newblock Cap'n {Proto}: An insanely fast data interchange format.
\newblock \url{https://capnproto.org}, 2013.

\bibitem[Chiusano and Bjarnason(2019)]{chiusano2019unison}
P.~Chiusano and R.~Bjarnason.
\newblock Unison: A new approach to distributed programming.
\newblock \url{https://www.unison-lang.org}, 2019.

\bibitem[Conway et~al.(2012)]{conway2012logic}
N.~Conway, W.~R. Marczak, P.~Alvaro, J.~M. Hellerstein, and D.~Maier.
\newblock Logic and lattices for distributed programming.
\newblock In \emph{Proc.\ ACM SoCC}, 2012.

\bibitem[Dennis and Van~Horn(1966)]{dennis1966programming}
J.~B. Dennis and E.~C. Van~Horn.
\newblock Programming semantics for multiprogrammed computations.
\newblock \emph{Communications of the ACM}, 9(3):143--155, 1966.

\bibitem[Epstein et~al.(2011)]{epstein2011towards}
J.~Epstein, A.~P. Black, and S.~Peyton~Jones.
\newblock Towards {Haskell} in the cloud.
\newblock In \emph{Proc.\ ACM Haskell Symposium}, pages 118--129, 2011.

\bibitem[Fischer et~al.(1985)]{fischer1985impossibility}
M.~J. Fischer, N.~A. Lynch, and M.~S. Paterson.
\newblock Impossibility of distributed consensus with one faulty process.
\newblock \emph{Journal of the ACM}, 32(2):374--382, 1985.

\bibitem[Gilbert and Lynch(2002)]{gilbert2002brewer}
S.~Gilbert and N.~Lynch.
\newblock Brewer's conjecture and the feasibility of consistent, available,
  partition-tolerant web services.
\newblock \emph{ACM SIGACT News}, 33(2):51--59, 2002.

\bibitem[Google(2008)]{protobuf2008}
Google.
\newblock Protocol {Buffers}: Language-neutral, platform-neutral extensible
  mechanisms for serializing structured data.
\newblock \url{https://protobuf.dev}, 2008.

\bibitem[Hayashibara et~al.(2004)]{hayashibara2004phi}
N.~Hayashibara, X.~D{\'e}fago, R.~Yared, and T.~Katayama.
\newblock The $\varphi$ accrual failure detector.
\newblock In \emph{Proc.\ IEEE SRDS}, pages 66--78, 2004.

\bibitem[Hellerstein(2010)]{hellerstein2010declarative}
J.~M. Hellerstein.
\newblock The declarative imperative: Experiences and conjectures in
  distributed logic.
\newblock \emph{ACM SIGMOD Record}, 39(1):5--19, 2010.

\bibitem[Hewitt et~al.(1973)]{hewitt1973actors}
C.~Hewitt, P.~Bishop, and R.~Steiger.
\newblock A universal modular {ACTOR} formalism for artificial intelligence.
\newblock In \emph{Proc.\ IJCAI}, pages 235--245, 1973.

\bibitem[Jacobs(1999)]{jacobs1999categorical}
B.~Jacobs.
\newblock \emph{Categorical Logic and Type Theory}.
\newblock Studies in Logic and the Foundations of Mathematics.
  Elsevier, 1999.

\bibitem[Lamport(1978)]{lamport1978time}
L.~Lamport.
\newblock Time, clocks, and the ordering of events in a distributed system.
\newblock \emph{Communications of the ACM}, 21(7):558--565, 1978.

\bibitem[Lamport(1998)]{lamport1998paxos}
L.~Lamport.
\newblock The part-time parliament.
\newblock \emph{ACM Transactions on Computer Systems}, 16(2):133--169, 1998.

\bibitem[Liskov(1988)]{liskov1988promises}
B.~Liskov and L.~Shrira.
\newblock Promises: Linguistic support for efficient asynchronous procedure
  calls in distributed systems.
\newblock In \emph{Proc.\ ACM PLDI}, pages 260--267, 1988.

\bibitem[Miller(2006)]{miller2006robust}
M.~S. Miller.
\newblock \emph{Robust Composition: Towards a Unified Approach to Access
  Control and Concurrency Control}.
\newblock PhD thesis, Johns Hopkins University, 2006.

\bibitem[Miller et~al.(2003)]{miller2003capability}
M.~S. Miller, K.-P. Yee, and J.~Shapiro.
\newblock Capability myths demolished.
\newblock Technical Report SRL2003-02, Johns Hopkins University, 2003.

\bibitem[Milner(1999)]{milner1999communicating}
R.~Milner.
\newblock \emph{Communicating and Mobile Systems: the $\pi$-Calculus}.
\newblock Cambridge University Press, 1999.

\bibitem[Ongaro and Ousterhout(2014)]{ongaro2014raft}
D.~Ongaro and J.~Ousterhout.
\newblock In search of an understandable consensus algorithm.
\newblock In \emph{Proc.\ USENIX ATC}, pages 305--319, 2014.

\bibitem[Sangiorgi and Walker(2001)]{sangiorgi2001pi}
D.~Sangiorgi and D.~Walker.
\newblock \emph{The $\pi$-Calculus: A Theory of Mobile Processes}.
\newblock Cambridge University Press, 2001.

\bibitem[Shapiro et~al.(2011)]{shapiro2011comprehensive}
M.~Shapiro, N.~Pregui{\c{c}}a, C.~Baquero, and M.~Zawirski.
\newblock A comprehensive study of convergent and commutative replicated
  data types.
\newblock INRIA Research Report RR-7506, 2011.

\bibitem[Van~Renesse et~al.(1998)]{vanrenesse1998gossip}
R.~Van~Renesse, Y.~Minsky, and M.~Hayden.
\newblock A gossip-style failure detection service.
\newblock In \emph{Proc.\ IFIP Middleware}, pages 55--70, 1998.

\bibitem[Virding et~al.(1996)]{virding1996concurrent}
R.~Virding, C.~Wikstr{\"o}m, and M.~Williams.
\newblock \emph{Concurrent Programming in Erlang}.
\newblock Prentice Hall, 2nd edition, 1996.

\bibitem[Wadler(2015)]{wadler2015propositions}
P.~Wadler.
\newblock Propositions as sessions.
\newblock \emph{Journal of Functional Programming}, 24(2--3):384--418, 2015.

\bibitem[Waldo et~al.(1996)]{waldo1996note}
J.~Waldo, G.~Wyant, A.~Wollrath, and S.~Kendall.
\newblock A note on distributed computing.
\newblock Technical Report TR-94-29, Sun Microsystems Laboratories, 1994.
\newblock Published in \emph{Mobile Object Systems}, LNCS 1222, Springer, 1996.

\bibitem[Yuan et~al.(2014)]{yuan2014simple}
D.~Yuan, Y.~Luo, X.~Zhuang, G.~R. Rodrigues, X.~Zhao, Y.~Zhang,
  P.~U. Jain, and M.~Stumm.
\newblock Simple testing can prevent most critical failures: An analysis of
  production failures in distributed data-intensive systems.
\newblock In \emph{Proc.\ USENIX OSDI}, pages 249--265, 2014.

\end{thebibliography}

% ===========================================================================
% APPENDIX
% ===========================================================================
\appendix

\section{Full Message Frame Specification}
\label{app:frame}

\begin{table}[H]
\centering
\caption{Message Frame Fields}
\begin{tabularx}{\textwidth}{llX}
\toprule
\textbf{Field} & \textbf{Size} & \textbf{Description} \\
\midrule
Magic          & 2 bytes  & Protocol identifier (\texttt{0x4A50}, ``JP'') \\
Flags          & 1 byte   & Bit flags: compressed (0x01), encrypted (0x02),
                             priority (0x04), migration (0x08),
                             capability (0x10) \\
Message Type   & 4 bytes  & Type tag derived from the algebraic data type hash \\
Version        & 2 bytes  & Schema version number \\
Length         & 4 bytes  & Payload length in bytes \\
Source PID     & 16 bytes & Source process identifier (8 bytes node ID + 8 bytes
                             process ID) \\
Dest PID       & 16 bytes & Destination process identifier \\
Capability Token & 0 or 32 bytes & Optional capability token; present when the
                             \texttt{capability} flag (0x10) is set in Flags.
                             Contains a 32-byte cryptographically signed token
                             authorizing the operation. Omitted (zero bytes)
                             for ordinary messages that require no capability. \\
Payload        & variable & Serialized message bytes \\
\bottomrule
\end{tabularx}
\end{table}

\section{Serialization Encoding Details}
\label{app:encoding}

\subsection{Primitive Encodings}

\begin{table}[H]
\centering
\caption{Primitive Type Wire Encodings}
\begin{tabularx}{\textwidth}{llX}
\toprule
\textbf{Type} & \textbf{Encoding} & \textbf{Notes} \\
\midrule
\texttt{Bool}    & 1 byte (0x00 or 0x01)           & \\
\texttt{Int}     & LEB128 (variable length)         & Zigzag encoding for
                                                       signed integers \\
\texttt{Float}   & 8 bytes IEEE 754 big-endian      & \\
\texttt{Float32} & 4 bytes IEEE 754 big-endian      & \\
\texttt{Char}    & 4 bytes UTF-32                   & \\
\texttt{String}  & LEB128 length + UTF-8 bytes      & \\
\texttt{Bytes}   & LEB128 length + raw bytes        & \\
\texttt{Unit}    & 0 bytes                          & \\
\texttt{Pid}     & 16 bytes (node ID + process ID)  & \\
\bottomrule
\end{tabularx}
\end{table}

\subsection{Composite Encodings}

\begin{table}[H]
\centering
\caption{Composite Type Wire Encodings}
\begin{tabularx}{\textwidth}{lX}
\toprule
\textbf{Type} & \textbf{Encoding} \\
\midrule
Record $\{l_1:\tau_1, \ldots\}$ & For each field: LEB128 field tag +
  LEB128 field length + encoded value. Fields in tag order. Terminated
  by tag 0. \\
\addlinespace
Variant $C_i(\tau)$ & LEB128 variant tag + encoded payload. \\
\addlinespace
\texttt{List[$\tau$]} & LEB128 count + encoded elements in order. \\
\addlinespace
\texttt{Map[$k$,$v$]} & LEB128 count + (encoded key + encoded value)
  pairs in key order. \\
\addlinespace
\texttt{Option[$\tau$]} & 1 byte tag (0x00 = None, 0x01 = Some) +
  encoded value if Some. \\
\addlinespace
\texttt{Result[$a$,$e$]} & 1 byte tag (0x00 = Ok, 0x01 = Err) +
  encoded value. \\
\bottomrule
\end{tabularx}
\end{table}

\section{Node Discovery Protocol Messages}
\label{app:discovery}

\begin{lstlisting}
type DiscoveryMsg =
  | JoinRequest({
      name: String,
      listen_addr: String,
      capabilities: List<NodeCapability>,
      cookie_hash: Bytes,
    })
  | JoinAccepted({
      members: List<NodeInfo>,
      cluster_state: ClusterState,
    })
  | JoinRejected(JoinError)
  | GossipPush({
      membership: MembershipState,
      vector_clock: VectorClock,
    })
  | GossipPull({
      membership: MembershipState,
      vector_clock: VectorClock,
    })
  | Heartbeat({
      from: NodeId,
      timestamp: Timestamp,
      load: Float,
    })
  | HeartbeatAck({
      from: NodeId,
      timestamp: Timestamp,
    })
\end{lstlisting}

\section{Phi Accrual Failure Detector Implementation}
\label{app:phi}

\begin{lstlisting}
type PhiDetector = {
  window_size: Int,
  threshold: Float,
  heartbeat_history: CircularBuffer<Timestamp>,
}

fn phi_value(detector: PhiDetector, now: Timestamp) -> Float =
  let intervals = compute_intervals(detector.heartbeat_history)
  let mean = Stats.mean(intervals)
  let std_dev = Stats.std_dev(intervals)
  let time_since_last = now - CircularBuffer.last(detector.heartbeat_history)
  -- Phi = -log10(1 - CDF_normal(time_since_last, mean, std_dev))
  let p = Stats.normal_cdf(Float.from_int(time_since_last), mean, std_dev)
  if p >= 1.0 then 100.0  -- cap at high value
  else -Float.log10(1.0 - p)

fn is_unreachable(detector: PhiDetector, now: Timestamp) -> Bool =
  phi_value(detector, now) > detector.threshold
\end{lstlisting}

\end{document}
